% Удалены все голые цифры в начале!
\usepackage{indentfirst}
\setlength{\parindent}{1.25cm}

\usepackage{enumitem}
\setlist{
	leftmargin=0pt,
	nosep,
	wide=\parindent,
	labelsep=0.5em,
	listparindent=\parindent
}
\setlist[itemize,1]{label=--} 
\setlist[enumerate,1]{label=\arabic*.}

\usepackage{titlesec}
\newcommand{\nohyphens}{\hyphenpenalty=10000\exhyphenpenalty=10000\sloppy}

\titleformat{\chapter}[hang]{\nohyphens\bfseries}{\thechapter}{1em}{}
\titleformat{\section}{\nohyphens\bfseries}{\thesection}{1em}{}
\titleformat{\subsection}{\nohyphens\bfseries}{\thesubsection}{1em}{}

\titlespacing*{\chapter}{\parindent}{0pt}{1em}
\titlespacing*{\section}{\parindent}{1em}{1em}
\titlespacing*{\subsection}{\parindent}{1em}{1em}

\usepackage{secdot}
\sectiondot{section}
\sectiondot{subsection}

\titleformat{name=\chapter,numberless}[block]
{\nohyphens\bfseries\filcenter}
{}
{0pt}
{\MakeUppercase}
\titlespacing*{name=\chapter,numberless}{0pt}{0pt}{1em}

\usepackage{float}
\captionsetup{justification=centering, singlelinecheck=false}
\DeclareCaptionLabelFormat{simple}{#1~#2}
\captionsetup{labelformat=simple}

\urlstyle{same}

\DeclareCaptionLabelSeparator{emdash}{ --- }
\captionsetup[table]{
	labelsep=emdash,
	justification=raggedright,
	singlelinecheck=false,
	position=top,
	skip=6pt
}

\usepackage{listings}
\usepackage{xcolor}

% Определяем цвета для JSON
\definecolor{json_key}{RGB}{0,0,255}       % Синий для ключей
\definecolor{json_string}{RGB}{163,21,21} % Красный для значений
\definecolor{json_background}{RGB}{250,250,250} % Почти белый фон

\lstdefinelanguage{json}{
	basicstyle=\ttfamily\footnotesize, % Уменьшенный шрифт (footnotesize)
	numbers=none,
	stepnumber=1,
	numbersep=8pt,
	showstringspaces=false,
	breaklines=true,
	backgroundcolor=\color{json_background},
	xleftmargin=2.5cm,      % Отступ слева (отодвигаем от края)
	xrightmargin=2cm,       % Отступ справа
	frame=none,             % Без рамки (по СТУ СФУ)
	stringstyle=\color{json_string},
	literate=
	*{0}{{{\color{json_key}0}}}{1}
	{1}{{{\color{json_key}1}}}{1}
	{2}{{{\color{json_key}2}}}{1}
	{3}{{{\color{json_key}3}}}{1}
	{4}{{{\color{json_key}4}}}{1}
	{5}{{{\color{json_key}5}}}{1}
	{6}{{{\color{json_key}6}}}{1}
	{7}{{{\color{json_key}7}}}{1}
	{8}{{{\color{json_key}8}}}{1}
	{9}{{{\color{json_key}9}}}{1}
	{:}{{{\color{black}{:}}}}{1}
	{,}{{{\color{black}{,}}}}{1}
	{\{}{{{\color{black}{\{}}}}{1}
	{\}}{{{\color{black}{\}}}}}{1}
	{[}{{{\color{black}{[}}}}{1}
	{]}{{{\color{black}{]}}}}{1},
	morestring=[b]",
	morestring=[d]'
}

% Цвета для подсветки
\definecolor{code_keyword}{RGB}{0,0,255}    % Синий для ключевых слов
\definecolor{code_string}{RGB}{163,21,21}   % Красный для строк
\definecolor{code_comment}{RGB}{0,128,0}    % Зеленый для комментариев
\definecolor{code_type}{RGB}{38,127,153}    % Бирюзовый для типов

\lstdefinelanguage{TypeScript}{
	keywords={export, async, function, return, method, body, await, string, const, let},
	keywordstyle=\color{code_keyword}\bfseries,
	ndkeywords={DocumentPayload, string},
	ndkeywordstyle=\color{code_type}\bfseries,
	identifierstyle=\color{black},
	sensitive=false,
	comment=[l]{//},
	morecomment=[s]{/*}{*/},
	commentstyle=\color{code_comment}\ttfamily,
	stringstyle=\color{code_string}\ttfamily,
	morestring=[b]',
	morestring=[b]",
	morestring=[b]{`} % Исправлено: теперь обратные кавычки (backticks) работают корректно
}

\usetikzlibrary{positioning, fit, shapes.geometric, arrows.meta, calc, backgrounds}

\setlength{\bibhang}{0pt}

% 1. Обнуляем стандартный висячий отступ biblatex
\setlength{\bibhang}{0pt}

% 2. Переопределяем окружение библиографии под требования СФУ (как wide list)
\defbibenvironment{bibliography}
{\list
	{\printtext[labelnumberwidth]{%
			\printfield{labelprefix}%
			\printfield{labelnumber}}
	}
	{
		\setlength{\leftmargin}{0pt}    % Вторая строка и далее — от левого края (0см)
		\setlength{\itemindent}{1.75cm} % Первый символ (номер) — отступ 1.25см
		\setlength{\labelsep}{0.5em}    % Расстояние от номера до текста
		\setlength{\labelwidth}{0pt}
		\setlength{\itemsep}{\bibitemsep}
		\setlength{\parsep}{\bibparsep}
	}
}
{\endlist}
{\item}

\DefineBibliographyStrings{russian}{
	urlseen = {дата обращения:},
}
